\documentclass{frama-c-book}

\usepackage{ifthen}

\input{fc-macros.tex}

\newcommand{\mthread}{\plugin{Mthread}}
\newcommand{\html}{\lang{HTML}}
\newcommand{\pthreads}{\library{pthreads}}

% this is for compilation purposes
\def\evadir{../../src/plugins/eva}
\def\sharedir{\evadir/share/mthread}
\def\testdir{\evadir/tests/mthread}
\def\mthreadtexdir{./mthread}
\def\examplesdir{./examples/mthread}

\def\etc{\emph{etc}\ldots}
\def\eg.{\emph{eg}.}
\def\ie.{\emph{ie}.}
\def\resp.{\emph{resp}.}

\def\mysec#1{\S{\ref{sec:#1}}}
\def\mychap#1{Chapter~{\ref{chap:#1}}}
\def\myfig#1{Figure~{\ref{fig:#1}}}
\def\myappdx#1{Appendix~{\ref{sec:#1}}}

\newcommand{\framacversion}%
           {\input{../../VERSION} (\input{../../VERSION_CODENAME}\unskip)}

\title{The Mthread plug-in}
\author{Boris Yakobowski \\
with Richard Bonichon, David Bühler and Basile Desloges}

\fcversion{\framacversion}

\cclicence{by-sa}
\copyrightstarts{2011}

%==============================================================================
\begin{document}
\sloppy
\emergencystretch 3em

\maketitle
\tableofcontents


\chapter*{Foreword}
\markright{}
\addcontentsline{toc}{chapter}{Foreword}

This is the documentation of the \FramaC plug-in \mthread.

The content of this document corresponds to the version \framacversion of \FramaC.
However the development of the \mthread plug-in is still ongoing:
features described here may evolve in the future.

\section*{Acknowledgements}

\medskip

\acknowledgeANR{the ANR project Veridyc (ANR-09-SEGI-016)}


\chapter{An introduction to \mthread}

\section{What is \mthread?}
\label{sec:mthread}

\mthread is a \FramaC plug-in dedicated to the analysis of concurrent
\C programs. It finds and displays multithreaded events,
such as thread creation, mutex
locking, access to shared variables, \etc. \mthread then gives a very
simplified view of the source code, in which only source statements
relevant to the concurrent behavior of the program are left. It also
displays variables that are shared between threads, as well as data sent by
threads on messages queues. For each shared memory zone, the mutexes
that may protect it are automatically inferred, and possible race conditions
are reported.


\section{How \mthread works}
\label{sec:mt_desc}

\mthread performs sound and precise analyzes of concurrent
programs. It is built on top of the \Eva analysis of \FramaC, and uses
the latter to derive sound values (hence sometimes over-approximated) for
all the variables of the program, including those that are shared
between multiple threads. Schematically, \mthread's behavior can be summarized as
follows:
\begin{itemize}
\item Do a symbolic execution of the main thread; find the threads
  it launches.

\item Do a symbolic execution of the new threads, possibly discovering
  other new threads, which are then also executed symbolically.

\item From each thread, compute the set of variables it reads
  and writes, as well as the messages it tries to receive and send.

\item Compute the shared variables of the program, by detecting
  variables that are accessed concurrently (\ie. by at least two
  threads that are live at the same time). On such concurrent
  accesses, record which mutexes are being hold by the various threads.

\item Restart the whole process, reinjecting the results obtained so far:
  \begin{itemize}
  \item threads receiving messages from a message queue are given the
    values sent to this queue by the other threads;

  \item threads reading shared variables ``see'' the values they write
    in those variables, but also those written by the other threads.

  \end{itemize}

\item Iterate the process above until all threads agree on the
  information sent and exchanged during the execution of the program.
\end{itemize}

Reaching  a fixpoint of the above process means that a sound approximation
of the behavior of the program has been obtained, by construction.
More details on how \mthread works are given in \mychap{mthread-theory}.


\section{Running \mthread}

\mthread is a \FramaC plug-in, and is activated when launching \FramaC
by turning on the {\tt -mthread} switch at your shell prompt, as follows:
\begin{shell}
% frama-c <C files> -mthread <mthread options>
\end{shell}

The various options that configure the behavior of \mthread will be
given throughout this document, and are summarized in
\myappdx{options}. In order to be really useful,
\mthread however requires you to instrument the thread library used
in your \C project, as explained in \mysec{instrumentation}.


\chapter{\mthread theory}
\label{chap:mthread-theory}

The schema given in \S\ref{sec:mt_desc} already gives a faithful
representation of \mthread fixpoint-based approach.  Through the \Eva
analysis, we obtain information about a thread; we then reinject those
information into (future) analyzes of the other threads. Reaching a
fixpoint guarantees that all threads agree on the concurrent part of
the program, and that we have found an over-approximation of their
behavior. The sections below detail some of the computations \mthread
does during the iterations.


\section{Calling contexts}

Although the approach outlined above is simple, obtaining precise
results is not. Indeed, we must be careful not to compute too general
behaviors for the various threads, the risk being to get unusable
results. % For example,
% the \inout plugin of \FramaC is able to compute memory zones
% accessed by a program. However,
%
\mthread uses some callbacks made available by the \Eva analysis to
record the state of each function at the end of its evaluation.
In order to avoid losing precision, \mthread fuses
those states only when the \emph{calling contexts} are the same.
%
% In order to achieve this result, given a function \verb+f+, \mthread
% carefully distinguishes between two different calls to \verb+f+ if
% they occur in two different \emph{calling contexts}.
Formally, a calling context is the callstack that lead to
the execution of \verb+f+, taking the statements at which the
calls originated into account.
\begin{lstlisting}[numbers=left]
main () {
  g ();
}
void g () {
  int x, y;
  f(&x, 1);
  f(&y, 2);
}
void f (int *p, int v) {
  *p = v;
}
\end{lstlisting}
In the example above, there are two distinct calling contexts for the
function \verb+f+, namely \verb+<main, 2>:<g, 6>+ and
\verb+<main, 2>:<g, 7>+.  By making a distinction between those two
calls, \mthread is able to detect that \verb+x+ (\resp. \verb+y+) is
always affected the value 1 (\resp. 2). This is much more precise than
the information available by only inspecting the state of the \Eva
analysis at the end of the execution, which merges together all the calls
to a function. (This can be easily verified by querying the
possible values for \verb+p+ and \verb+v+ in the GUI of \FramaC, which
would lead to conclude that \verb+x+ and \verb+y+ are affected
either 1 or 2, without the possibility to know which one).


\section{Concurrent control-flow graphs}
\label{sec:mthread-cfg}

\subsection{General idea}

One result of the analyzes done by \mthread is the \emph{concurrent
  control-flow graph} of each thread. Those graphs aim at
displaying all the following \emph{events}:

\begin{itemize}
\item calls to a \verb+mthread.h+ primitive;
\item accesses to a shared memory zone (see \S\ref{sec:shared-vars}).
\end{itemize}

Basically, we build a very high-level view of the
functions called by a thread, with the following characteristics:
\begin{enumerate}
\item only function calls that contain an event, or that transitively
  lead (through another call) to such an event, appear in the graph;

\item functions are duplicated for each calling context they appear in;

\item inside the body of a function, only events
  and high-level control-flow statements such as \verb+if+ and
  \verb+loop+ appear. Control-flow statements that do not lead to
  an event are also removed.
\end{enumerate}
Points~1 and~3 guarantee that the graph keeps a reasonable size,
even with very big programs. Indeed, most of the code is typically
not related to its concurrent structure. Conversely, point~2 expands
the size of the graph, but increases its precision. Indeed, the
statements executed by a function can be very different from one call
to another, and this is captured by our use of calling contexts.

Notice that the concurrent control-flow graph of a thread is very
different from what would be obtained with the \slicing plugin of
\FramaC.  In particular, our control-flow graph does not represent
executable code at all. (For example, incrementations of loop indices
are generally removed from the graph.) Conversely, our graph can be more
precise when a function is called multiple times, and roughly corresponds
to the specialization obtained by \verb+-slicing-level 3+.

\subsection{Example}

\begin{figure}[htb]
  \centering
    \includegraphics[scale=0.4]{\examplesdir/ccfg.pdf}
  \caption{Concurrent control-flow graph for the main thread of our example}
  \label{fig:ccfg}
\end{figure}

The concurrent control-flow graph for the main thread of the
example \verb+doc/eva/examples/mthread/ccfg.c+, reprinted below, is given in
Figure~\ref{fig:ccfg}. How to generate this graph is explained in
\S\ref{sec:html}.

\lstinputlisting[numbers=left,language={[ANSI]C}]{\examplesdir/ccfg.c}

Let us illustrate through our example the characteristics of concurrent
control-flow graphs that have been mentioned above. The options we hint at
are documented in \myappdx{options}.
\begin{itemize}
\item The topmost node contains the name of the function
  the thread starts with, here \verb+main+.

\item Calls to other functions are inlined within the graph, as can be seen
  \eg. for \verb+g1+. A dotted grey edge links \verb+Call+ nodes to
  the corresponding \verb+return+ ones. (Option \verb+-return-edges+.)

\item Functions called twice in two different calling contexts,
  \eg. \verb+g1+, are inlined twice. Each body shown represents
  precisely the execution of the corresponding call. For \verb+g1+,
  the first call creates the first thread, while the second call has a
  behavior that depends on \verb+i+.

\item For function calls occurring through pointers \eg. the second
  call to \verb+g1+, the real name of the function is printed between
  the \verb+Call+ node and the body of the function.

\item Calls to functions that do not lead to an event are removed. For
  example, the call \verb+g2(arr)+ does not appear anywhere.

\item Nodes with a red border represent immediate calls to one or more
  \verb+mthread.h+ primitives, that are listed in the node. (In this
  case, only thread creations occur.)

\item Nodes with a blue or green background represent accesses to
  a shared zone of the memory, and are discussed in \S\ref{sec:shared-vars}.

\item \verb+Loop+ nodes represent \verb+while(1)+ loops;
  \verb+for+ loops are automatically desugared into \verb+while+
  ones by \FramaC.

\item Diamond nodes appear for a function without
  definition, \eg. \verb+g2+.  They represent all the events that occur
  during the call to the function inside a single node. Functions
  without definition use their  ACSL prototype to specify the data
  they read and write.

\item \verb+if+ constructs for which the condition is either
  completely true or completely false in the given context are
  removed. This is the case for the \verb+if (!x)+, as \verb+x+ is
  always equal to 0 in the program. Similarly, loops whose body do not
  contain an event are removed. Those simplifications can be
  deactivated with \verb+-mt-full-cfg+.

\item The \verb+exit+ node represent the end of the thread, hence
  the outgoing edge that goes nowhere.

\item Although this is not shown here, the \mthread graph simplifier
  is sometimes forced to leave some nodes that do not really
  contribute to the concurrent structure of the program. This is
  typically the case for functions that use \verb+goto+s. Those nodes
  will not have any border or background.

\end{itemize}


\subsection{Understanding loops in concurrent control flow graphs}

A word must be said on the composite node \verb|Create thread jobs[1..3]|
of our example. It must \emph{not} be understood as ``at each iteration
of the loop, three threads are created''. This is indeed impossible:
\begin{itemize}
\item each node corresponds to a single statement,
  and no \verb+mthread.h+ primitive can create three threads in a single
  statement;
\item \mthread does not allow the same thread to be launched  more than once,
  as indicated in \mysec{mthread-warnings}.
\end{itemize}

The correct way to read the graph is the following: at each iteration
of the loop, a different thread is created, with some iterations
possibly spawning none. (In fact, in our case the
iteration for $i=4$ does not create a thread.)


\section{Shared zones}
\begingroup
\label{sec:shared-vars}

\def \ifempty #1#2#3{\ifthenelse{\equal{#1}{}}{#2}{#3}}
\def\ri#1{\mathsf{Ri}(#1)}
\def\wi#1{\mathsf{Wi}(#1)}
\def\Zi{\mathsf{RWi}}
\def\Wp#1#2{\mathsf{Wp}^{#1}(#2)}
\def\Rp#1#2{\mathsf{Rp}^{#1}(#2)}
\def\live#1#2{\mathsf{live}(#1@#2)}
\def\shared#1{\mathsf{shared}\ifempty{#1}{}{(#1)}}


In this section, we call \emph{shared zone} a region of the memory on
which a race condition between at least two threads can occur.
%
\mthread performs a fine-grained analysis to detect those
regions. % A simple solution to do so
% would be to reuse \inout plugin of \FramaC. By calling it on a
% thread, we would obtain the list of globals it accesses.  Then it
% would suffice to intersect the results for all distinct threads,
% and keep all memory zones that are in any intersection.
It proceeds as follows:
\begin{enumerate}
\item Once a thread is evaluated by the \Eva analysis, we compute the
  global variables it reads and writes, using the \inout plugin
  of \FramaC. This plugin uses the results of the \Eva analysis,
  thus giving us a sound but quite imprecise over-approximation
  of the shared zones accessed by this thread. Let us call
  $\ri{j}$ (\resp. $\wi{j}$) the zones that are read
  (\resp. written) by the thread $j$.

\item After a full iteration (once all threads have been computed),
  we find all the zones that are read by at least one thread and written by
  at least one another.
  $$\Zi=\bigcup_{j,k,j\neq k}(\ri{j} \cap \wi{k})$$
  This set over-approximates the shared zones, and we
  call it \emph{potential shared zones}.

\item For each thread that accesses a zone in $\Zi$
  we start another \Eva analysis, and \emph{watch}
  the zones above: at the end of the execution of any
  function, if it reads or writes a zone in $\Zi$,
  we record a \mthread event for this access. This
  event will thus appear in the control-flow graph for the current
  thread.

  Let $\Rp z j$ (\resp. $\Wp z j$) be the set of those precise
  read (\resp. write) events relative to the zone  $z$, for the thread $j$.

\item Once all the needed threads have been recomputed, we compute
  the threads that are live on each point of the control-flow graphs.
  Let us note $\live j e$ the fact that the thread $j$
  is live at the node containing the event $e$.

  By definition, there is potentially a race condition on a zone $z$
  if it is written by one thread and read by another, both threads being
  live at the same time.
  %
  Thus,  for each zone z of $\Zi$, we define the fact it is shared by:
%
  $$\shared z=\exists {j,k},\ j\neq k \land \big(
       \exists e_j \in \Rp z j,\
        \exists e_k \in \Wp z k,\
       \live j {e_k} \land \live k {e_j}\big)
  $$

\end{enumerate}
(Of course, this is only the mathematical definition of the
$\shared{}$ predicate. The computations themselves are done
efficiently, to avoid the cubic complexity of the formulas
above.) % Moreover, we do not recompute the threads each time
% $\Zi$ changes. Instead, it is needed

\smallskip The definition of $\shared{}$ is as precise as possible given
the information available to \mthread,
while remaining sound. In particular, it avoids flagging as shared an
important set of variables, those that are only initialized (\ie. written)
by the main thread, and used (\ie. read) afterward by the various threads.
As long as the initialization occurs before the creation of any
of the threads that access the variable, this variable is \emph{not}
shared.\footnote{In fact, to reduce the time spent computing shared
zones, \mthread completely ignores all the accesses that occur before
the creation of the first thread.}

\smallskip

Once all the analyzes are finished, \mthread classifies the events
representing accesses to potential shared zones in three categories.
Let us consider an event $e$ for an access to a zone $z$.
\begin{description}

\item [Non-shared access]

  This means that $z$ is in fact \emph{not} a shared zone.
  Although $z$ was in $\Zi$, there is never any race condition
  when accessing this zone. Since those zones are not
  important, $e$ is not shown in the control-flow graph by default.
  This can be overridden by the option \verb+-mt-non-shared-accesses+
  if desired.

\item [Shared, non-concurrent access]

  \mthread has determined that $z$ is indeed a shared zone. However,
  the particular event represented by $e$ is not concurrent, because
  all the other threads that access $z$ are either not created yet, or
  canceled. This is typically the case for most initializations of
  shared zones by the main thread. In the control-flow graph,
  those events are shown in green. The option
  \verb+-mt-no-non-concurrent-accesses+ can be used to hide them if desired.

\item [Concurrent access]

  The zone $z$ is indeed a shared zone, and the access is concurrent.
  That is, a race condition is possible at $e$. It is shown in
  blue in the control-flow graph.
\end{description}
%
An example of the various cases above can be found in the file
\verb+src/plugins/eva/tests/mthread/sharedvars.c+,
which we do not reproduce below for space consideration.
Of the 6 variables of the programs, 3 are shared (those starting by \verb+'s'+)
and 3 are not (those starting by \verb+'u'+).
Running \mthread on it with the option \verb+-mt-verbose 3+ is concluded by
%
% Taken from tests/sharedvars.res.oracle
\begin{logs}
[mt] Imprecise locations to watch: u3; s4; s5; s6
[mt] Possible read/write data races:
  s6:
    read by jobs4 at sharedvars.c:56, unprotected
    read by jobs6 at sharedvars.c:75, unprotected
    write by jobs4 at sharedvars.c:57, unprotected
    write by jobs6 at sharedvars.c:76, unprotected
  s5:
    read by jobs51 at sharedvars.c:62, unprotected
    write by jobs51 at sharedvars.c:63, unprotected
    write by jobs5 at sharedvars.c:69, unprotected
  s4:
    read by jobs4 at sharedvars.c:53, unprotected
    write by <main> at sharedvars.c:100, unprotected
    write by jobs4 at sharedvars.c:54, unprotected
[mt] Shared memory: s4; s5; s6
\end{logs}
The line ``Imprecise locations to watch'' indicates that the
potential shared zones are the variables \verb+u3+, \verb+s4+,
\verb+s5+ and \verb+s6+ . The section ``Possible read/write data races'' and
the line ``Shared memory'' however indicates that \verb+u3+ is not
really shared.

Also, examining the control-flow graph of the main thread shows a
non-concurrent access to \verb+s4+ before the creation of
\verb+&jobs4+.  This access is not listed above, as it is not
concurrent ---and thus must not be taken into account when examining
the mutexes that protect \verb+s4+.

\subsection{Protecting shared zones through mutexes}

The race conditions evoked in the previous section are theoretical.
That is, they can be prevented using an appropriate use of mutexes.
However, once all the shared zones have been found, \mthread needs
to do very little more to have this information.

Indeed, for each access to a shared zone (\ie. an event in the
control-flow graph), we known which mutexes are locked, and which are
not. Thus, in its final output, when \mthread lists all the shared
zones it has detected, it adds the information it possesses about
mutexes. Mutexes that are guaranteed to be locked are written
directly.  Mutexes that may or may not be locked (\eg. that are locked
in one branch of the program, but not in another) are prefixed by
\verb+(?)+.

In a second time, \mthread combines those information together and
list for each zone the mutexes that are either possibly or
systematically locked when the zone is accessed. A shared zone that is
protected by at least one guaranteed mutex will not be subject to a race
condition.

Since \verb+sharedvars.c+ does not use mutexes at all, it is not very
pertinent here. Some examples of protection outputs are given in
\S\ref{sec:results-philo}.


\endgroup

\section{Related works}

Ferrara \cite{Ferr09} uses a fixpoint-based approach very similar to
ours to analyze Java bytecode.
%
The static analyzer \tool{Locksmith} \cite{hicks06lockinference}, which
is dedicated to finding data races in multithreaded \C programs, possess
some similarities with our shared zones detection
algorithm.
%
The \tool{Goblint} \cite{VVojVVen07} is race-detection tool using some
fixpoint computation (resolved by a constraint solver): it offers a
path-sensitive analysis of data-races, based upon conditional
constraint propagation and points-to analysis.
%
Miné~\cite{mine-LMCS12} builds an analyzer for concurrent code on
top of the Astree abstract interpreter. Apart from the use of two
distinct base analysers, our approach and his are very similar.


\chapter{Instrumenting the \C concurrent primitives}
\label{sec:instrumentation}

To precisely detect calls to concurrent primitives during the
symbolic execution of the program, \mthread makes the hypothesis
that those primitives invoke low-level \mthread functions. Hence,
the first step in using \mthread consists in properly stubbing
the thread library of the program. This work has already been done
for parts of the \pthreads library.
The functions currently supported by
\mthread are detailed in the next sections.

In this chapter, and unless stated otherwise, the files we refer to
are located in the \verb+src/plugins/eva/share/mthread/+ directory of \FramaC
sources, or alternatively in the \verb+share/mt+ of the Frama-C installation.
After installation, this directory can be found by running the command
\verb+echo $(frama-c -print-share-path)/mt+. This directory will be
abbreviated as \verb+$MTSHARE+ in the rest of this document.

\section{First steps}

For a new project, the first step consists in finding within
the \C sources the {\sf .h} file containing the declarations
for the various multithreaded primitives used in the code.
In general, those functions include at least
\begin{itemize}
\item thread creation (and possibly cancellation);
\item mutex locking and release;
\item emission and reception on/from a message queue.
\end{itemize}
Other interesting primitives are those initializing the
structures used to refer to the objects above (threads, mutexes,
queues), functions using more evolved concurrency primitives
(spinlocks,\ldots) \etc. A detailed status of which functions
are currently handled by \mthread is given in
\myappdx{stubbed-functions}.

Once the prototypes of the functions above have been found, any
potential implementation must be removed from the source, for example
by using well-placed \verb+#ifdef 0+ lines. This step is however
typically not needed, as these functions usually belong to the OS
implementation, the source code of which is rarely available.



\section{Stubbing the header (\texttt{.h}) files}

The next steps consists in stubbing the existing concurrency \verb+.h+ files.
The primary responsibility of this step is to define all the types
in the prototypes of the functions in terms of either
\verb+__fc_mthread_id+ or \verb+__fc_mthread_name+. Both types are
defined in the \mthread header \verb+mthread.h+ as

\begin{lstlisting}[language={[ANSI]C}]
typedef void *__fc_mthread_name;
typedef int __fc_mthread_id;
\end{lstlisting}

In general, \verb+__fc_mthread_id+ is the return type of the initialization
functions, and also the type used by functions that {\it use} an object.
During execution, they are simply sequential non-null offsets to an array
allocated by \mthread, that itself holds the state of the object. The non-null
information is important, as some code uses the convention \verb+ v == 0+ to
test whether an object is initialized. Also, we cannot return a pointer, as
some concurrenct library assume that thread ids are no bigger than the
\verb+short+
type; returning short integers (unless the code allocates an inordinate
amount of \eg. mutexes) ensures that our ids can safely be cast to
\verb+short+, or even \verb+char+.

By contrast, \verb+__fc_mthread_name+ is the input type used by
initialization functions. It is used as a hint to name the mutex, thread,
or queue. It can be either \verb+NULL+, in which case \mthread will
use an internal name, a constant string, or the address of a global variable,
with possibly an offset (if the variable is an array cell).

As an example, let us show how those two types are used in the
prototypes of the primitive \mthread functions. The lines below are
also an excerpt of \verb+mthread.h+. (The entire file is given in
\myappdx{stubbed-functions}.)
\begin{lstlisting}[language={[ANSI]C}]
__fc_mthread_id Frama_C_thread_create(__fc_mthread_name,
                                         void *(*)(void *),
                                         ...);
int Frama_C_thread_cancel(__fc_mthread_id);

__fc_mthread_id Frama_C_mutex_init(__fc_mthread_name);
int Frama_C_mutex_lock(__fc_mthread_id);
\end{lstlisting}

\medskip

To conclude this section, let us consider excerpts of the stubbing that has been
done for queue-related functions. The prototypes can be found in the file
\verb+mthread_queue.h+, and are reprinted below.

%% CHECKRELEASE
\lstinputlisting[breaklines=true,firstline=12,lastline=18,language={[ANSI]C}]{\sharedir/mthread_queue.h}

Except for the \verb+typedef+ declaration, everything can be taken
verbatim from the original header that is being stubbed. The type
\verb+msgqueue_t+ is defined as type alias to \verb+__fc_mthread_id+. Indeed,
the interface of the message queue library does not lend itself to the
separation between id and name we use in \mthread, and the same can be said for
\pthreads. Instead, as the next section will show, the initialization primitives
use the address of the object they create when naming them.

If the types are compatible it is possible to write the stubs for a library
without having to provide an updated header file. This is the case for instance
with \pthreads where the stubs are written to be compatible with the system
header file shipped with Frama-C's libc in the file \verb+mthread_pthread.c+,
as explained in the next section.


\section{Stubbing the source (\texttt{.c}) files}

The bulk of the stubbing work consists in implementing the concurrent
\C primitives in terms of the low-level \mthread ones.
Stubs are generally very easy to write, as most of the time they consist in:
\begin{itemize}
\item disregarding useless arguments (such as initialization options
  \mthread may not handle yet), or swapping some arguments around;
\item dereferencing pointers, if a pointer is supplied while
  \mthread needs the value it points to;
\item for initialization functions, storing or returning the result of the call
  to the low-level \mthread primitive (which will be the id of
  the thread, mutex or queue for \mthread);
\item translating the \mthread return codes into those of the OS library.
\end{itemize}


\subsection{\pthreads library}

The stubbing for the \pthreads
library can be found in \verb+mthread_pthread.c+. Its interesting
parts are reprinted below.

%% CHECKRELEASE
\lstinputlisting[breaklines=true,firstline=12,lastline=58,language={[ANSI]C}]{\sharedir/mthread_pthread.c}
Notice the recurring pattern \verb+*obj = obj_init(&obj, ...)+ (with proper
error-handling) used in
both functions \verb+pthread_create+ and \verb+pthread_mutex_init+.
The address holding the object is used as name hint during the
creation. % (it is casted to \verb+__fc_mthread_name+ by \FramaC parser)
Then the \mthread initialization function returns the id of the object, which
is stored at the given address. The functions that use this id
either dereference their argument if the id is passed as a pointer
(\verb+pthread_mutex_lock+) or use it directly otherwise
(\verb+pthread_cancel+), depending on their POSIX prototype.


As hinted by the comments, not all functions are properly stubbed. The
\verb+*setcancel+ functions, which are related to the cancellability of
a thread, are not given a body; for now, we implicitly assume that
\verb+pthread_cancel+ always succeed in stopping a
thread. For functions
that do not initialize values, there is little differences between
not stubbing a function, and giving it a trivial body; the latter
approach however silences a few \FramaC warnings.  For functions that
initialize a structure used later, a stub is however mandatory.
% (We have not yet stubbed functions such as \verb+pthread_attr_init+ for this
% reason.)

Also, thread return codes are not stored yet, which means that
\verb+pthread_join+ is not modeled as precisely as possible.


\chapter{Analyzing a full project without warnings}

This chapter explains the warning or error messages emitted by
\mthread during its analysis.  \mthread's own analysis can only be as
precise as the one done by \Eva for each thread. Thus, setting up
the latter correctly is important; relevant information can be found
in its own manual:
\begin{center}
\url{https://frama-c.com/eva.html}
\end{center}

During the analysis, the first hint that something might have gone
awry resides in the warnings sent back to the user. As a general rule
of thumb, it is good to eliminate those messages.  How to read the
\emph{results} of \mthread will be explained in \mychap{results}.


\section{The philosophers example}
\label{sec:philo}

\def\cfile{philo.c}

In the remainder of this document, we will use the source code below
to exemplify some uses of \mthread. It is taken from the file
{\tt doc/eva/examples/mthread/\cfile}, and contains a modified version of the
classic philosophers problem.

\lstinputlisting[firstline=3,numbers=left,language={[ANSI]C}]{\examplesdir/\cfile}

This code presents some interesting challenges. First, the ids for the
threads and mutexes of the program are stored in two arrays,
\verb+jobs+ and \verb+locks+ respectively. Both arrays are initialized
through loops --- a challenge for any analyzer. Moreover, the
behavior of the various threads is governed by a unique function: the
only difference between them lies in the argument they initially
receive. Finally, the various threads write in the global variable
\verb+end2+, and send a partially initialized message on the queue
\verb+queue+. The termination of the main thread is influenced by
those two objects.


\subsection{A first try}
\label{sec:philo-no-slevel}

Let us start \mthread on this program. We need to start \FramaC on both
\verb+philo.c+ and on our stubbed \pthreads library. The option
\verb+-mt-threads-lib+ with value \verb+pthreads+ should be used to tell
\mthread that we need the stubbed \pthreads library. Additionally,
\verb+philo.c+ uses message queue functions stubbed separately in
\verb+$MTSHARE/mthread_queue.c+, so this file should be added to the list of files
to parse.
Finally, the \Eva analysis is by default very verbose, and it is a
good idea to partially silence it using \verb+-eva-verbose 0+.
Thus, a complete invocation of \mthread would be
%% CHECKRELEASE
\begin{shell}
% frama-c -mthread -mt-threads-lib pthreads \
   $MTSHARE/mthread_queue.c philo.c \
   -eva-verbose 0
\end{shell}
(\FramaC assumes the main function is called \verb+main+, which is the
case here)

\smallskip

While \mthread's output is also rapidly verbose (we will only reproduce
snippets below), it is quite apparent that something has gone awry.
Many lines start by \verb+philo.c:ll[mt] warning:+, where \verb+ll+ is
a line number. The prefix \verb+[mt]+ being a short-name for \mthread,
let us examine a few of those warning lines.

The first one is
%% CHECKRELEASE
% Taken from examples/mthread/philo.noslevel.log
\begin{logs}
[mt] philo.c:46: Initializing mutex #2
[mt] philo.c:46: Warning:
  Trying to initialize mutex #2, but unable to check its precise status
  (internal value {0; 1} should be 0).
\end{logs}

By contrast, the previous line was
%% CHECKRELEASE
% Taken from examples/mthread/philo.noslevel.log
\begin{logs}
[mt] philo.c:46: Initializing mutex locks[0].
\end{logs}

Since line 46 is a call to \verb+pthread_mutex_init+, it is clear
that the mutex initialization did not succeed in all cases. This is
of course due to the presence of the loop, which is executed symbolically,
but imprecisely. While the first iteration of the body proceeds as
expected (and initializes \verb+locks[0]+), the next ones do not. Instead,
during the analysis of the loop, the value \verb+i+ ranges over the sets
\verb+{0}+ \ldots \verb+{0, 1, 2, 3, 4}+. While in the first iteration, \Eva can
reduce the value of \verb+locks[i]+ to a single value, this is not the case for
the following iterations where the expression ranges over the locations
\verb|&locks[1 ... 4]|. \mthread cannot distinguish one mutex from the other and
uses a generic name \verb+#2+ for the set of locks. Moreover, \mthread cannot
know which mutex is really being initialized at each iteration and accordingly
warns that it is unable to check if the mutex has not already been initialized.

In this particular case, since the initialization is in a loop, increasing
the \verb#-eva-slevel# option of the \Eva analysis should help. Indeed this
option controls, among others, the symbolic execution of loops. By default,
no \verb#-eva-slevel# is used, and variables modified inside loops become
imprecise immediately.

Before increasing this option, let us consider the other warnings.
For line 19, we actually get two different warnings depending on which thread
is analyzed.

The first warning which occurs during the analysis of thread \verb+jobs[0]+ is:
%% CHECKRELEASE
% Taken from examples/mthread/philo.noslevel.log
\begin{logs}
[mt] philo.c:19: Warning: Trying to lock a possibly uninitialized mutex.
[mt] philo.c:19: Warning:
  Trying to lock mutex #2, but unable to check its precise status
  (internal value {0; 1} should be 1).
\end{logs}
This is the same problem as before: mutex \verb+#2+ is actually a set of mutexes
and \mthread is not able to know which of the concrete mutexes is being locked.

The second warning occurs during the analysis of thread \verb+#3+:
%% CHECKRELEASE
% Taken from examples/mthread/philo.noslevel.log
\begin{logs}
[mt] philo.c:19: Warning:
  During mutex lock: invalid mutex id. Non-integer value: {0; 1; 2} Ignoring.
\end{logs}
This is similar to the previous warning in that \mthread cannot know which mutex
is being locked, but this time this is due to the fact that the threads too have
been merged into an abstract thread representing several concrete threads. We
can see that by checking the logs for the line 51 that creates the threads:
%% CHECKRELEASE
% Taken from examples/mthread/philo.noslevel.log
\begin{logs}
[mt] philo.c:51: New thread: jobs[0], fun job, parent <main>, args {0}
[mt] philo.c:51: Start thread jobs[0]
[mt] philo.c:51: New thread: #3, fun job, parent <main>, args {0; 1}
[mt] philo.c:51: Start thread #3
[mt] philo.c:51: New context for #3, fun job, parent <main>, args {0; 1; 2}
[mt] philo.c:51:
  New context for #3, fun job, parent <main>, args {0; 1; 2; 3; 4}
[mt] philo.c:51: Thread #3, fun job, parent <main>, args {0; 1; 2; 3; 4}
\end{logs}
Here we see that after creating a thread for \verb+jobs[0]+, \Eva starts merging
the states for several threads and \mthread creates a second thread generically
called \verb+#3+. Instead of creating several new threads, the context of this
abstract thread is updated to support more arguments until the set
\verb+{0; 1; 2; 3; 4}+ is used. Since this argument is used to select the
mutexes to lock and unlock in function \verb+aux+, \mthread cannot know which
mutex is being used.


\subsection{Unrolling loops}
\label{sec:slevel}

Let us increase the \verb+slevel+ to separate the states of the loops in the
\Eva analysis. Here, it suffices to execute loops precisely 5 times. Thus, we
add option \verb+-eva-slevel 5+ to the command-line used above. In this
case, it is sufficient to make all warnings disappear.

Alternatively, it is also possible to \emph{syntactically} unroll loops.
Although it is seldom useful for the \Eva analysis, it may be for \mthread
itself. Let us consider the file
\verb+doc/eva/examples/mthread/init.c+ of \mthread, which is reproduced below.
%
\lstinputlisting[breaklines=true,numbers=left,language={[ANSI]C}]{\examplesdir/init.c}
%
This example try to initialize 12 mutexes, using three different
mechanisms. In the first loop, the mutexes are named using
a location in an array. In the second loop, third and fourth loops, they are
named using constant strings in increasingly large loops. The result of the
analysis of the main thread by \mthread are given below (with a proper
\verb+-eva-slevel+ 3):
%% CHECKRELEASE
% Taken from examples/mthread/init.log
\begin{logs}
[mt] *** Computing value analysis for main thread
[eva:experimental] Warning: The mthread domain is experimental.
[mt] New thread: <main>, fun main
[mt] init.c:20: Initializing mutex locks[0]
[mt] init.c:20: Initializing mutex locks[1]
[mt] init.c:20: Initializing mutex locks[2]
[mt] init.c:23: Initializing mutex mu1
[mt] init.c:23: Initializing mutex mu2
[mt] init.c:23: Initializing mutex mu3
[mt] init.c:28: Initializing mutex mu4
[mt] init.c:28: Initializing mutex mu5
[mt] init.c:28: Initializing mutex mu6
[mt] init.c:32: Initializing mutex mu7
[mt] init.c:32: Initializing mutex #11
[mt] init.c:32: Warning:
  Trying to initialize mutex #11, but unable to check its precise status
  (internal value {0; 1} should be 0).
[mt] init.c:36: Initializing mutex 0
[mt] init.c:36: Warning:
  Trying to initialize mutex 0, but it might already be unlocked
  (and initialized).
[eva] init.c:13: Warning: The mutex 0 is already registered.
[mt] *** First value analysis for main thread done.
\end{logs}
As can be seen, the first 9 initializations succeed without
problem. The mutexes created in the first loop are named after the array
passed as hint (and the index of the mutex in the array), while the
exact names contained in \verb+names+ are used for the mutexes 4 to 9 in the
second and third loops.

Problems however arise in the last loop, although it is very similar
to the third one. \mthread initializes the tenth mutex on the first
iteration, but complains in the second that it cannot determine if the mutex has
already been initialized. Indeed the number of iterations of the loop is greater
than the value given to \verb+-eva-slevel+ and \mthread is not able to separate
the different mutexes. A possible solution (other than increasing the
\verb+-eva-slevel+ value) is to syntactically unroll the loop, as was done for
the third one. The instruction \verb+//@ loop unfold N;+ instructs \FramaC to
unroll a loop \verb+N+ times. Afterward, \mthread sees some mutex
initializations at different statements, and accepts them. Internally however,
those statements point to the same initial line number, hence the messages for
the third loop in the log.

\section{Other \mthread warnings}
\label{sec:mthread-warnings}

In this section, we give a short survey of some the warnings emitted by
\mthread.  In a log output, those warnings contain the string
\verb+[mt]+.  Most of the time, the warnings are self-explanatory, and
they sometimes contain their own solution. Roughly speaking, they can
be partitioned in the categories below.
\begin{description}
\item[Erroneous or imprecise arguments.]
\mthread systematically sanitizes the arguments it receives from
the \verb+Frama_C_*+ functions defined in \verb+mthread.h+, and ignores the
entire call (with a warning) when it cannot give a sense to them.
%
We have already given an example in \mysec{philo-no-slevel} with an
imprecise name for a mutex initialization. Nearly identical warnings
are emitted with imprecise names or ids, for threads, mutexes or queues.

Other similar errors can include passing a function without a body to
the thread creation function
\verb+Frama_C_thread_create+, or too few arguments. The corresponding
messages are given below.\footnote{The messages are obtained by simple
modifications of our examples and stubs, not shown in this document.}
\begin{logs}
philo.c:51:[mt] warning: During thread creation: invalid thread function.
                Missing definition for function 'job'. Ignoring.
\end{logs}
\begin{logs}
philo.c:51:[mt] user error: When creating thread &jobs[0] from function
                job: too few arguments, 1 expected but 0 given. Ignoring.
\end{logs}

\item[Multiple creation of a unique thread.]

\mthread is quite tolerant when it encounters code that would initialize
again a mutex or a queue potentially already initialized:
\begin{logs}
philo.c:46:[mt] warning: Mutex &locks might be already initialized
\end{logs}

We cannot be as lenient for threads, however, to preserve the
correctness of our analyzes. Thus, when we detect a thread that seems
to be started twice, we immediately fail.  Of course, it remains
possible to launch two threads with exactly the same arguments, but
the program must use two different names.
\begin{logs}
philo.c:51:[mt] Thread &jobs[0], fun job, parent _main_, args {0; }
philo.c:51:[mt] user error: Thread &jobs[0] has already been created
                previously in the current thread.
\end{logs}



\item[Read or write of the entire memory.]  If the \Eva analysis
  dereferences a very imprecise pointer, it can access the whole
  memory.  This completely invalidates the assumptions made by
  \mthread when it searches for shared memory, and can make it very
  imprecise. We therefore entirely ignore the access. Since such an
  imprecise pointer almost always comes from an erroneous stubbing, or
  a very buggy original code, this is not a limitation in practice.

  The directory \verb+tests+ contains an example designed to test this
  case, called \verb+read_all.c+. The \Eva analysis
  prints a warning when the faulty pointer \verb+p+ is being dereferenced
  (line~4). The \mthread warning is on line~7.

  \begin{logs}[numbers=left]
read_all.c:26:[mt] New thread: &jobs, fun f, parent <main>,
                   args [0..4294967295]
read_all.c:26:[mt] Start thread &jobs
read_all.c:28:[mt] user error: read of the whole memory. Ignoring to allow
     Mthread to continue, but the  analysis will not be correct.
read_all.c:28:[eva] warning: Completely invalid destination for assigns
      clause *p. Ignoring.
  \end{logs}

\item[Buffer overflow in message sending or receiving.]
\mthread send and receive functions take as input either a source
buffer, or a destination one, as well as its size. Of course, the
\Eva analysis must be wary of buffer overflow. Let us successively
change the declaration of the buffers \verb+buf+ and \verb+end+ of
\verb+philo.c+ to \verb+char[1]+ (lines 31 and 42).
%

\subsubsection{Small buffer during emission}
\begin{logs}
../share/mthread_queue.c:10:[kernel] warning: out of bounds read.
                assert \valid(mess+(0..size-1));
philo.c:25:[mt] Sending message on &queue, content [0..1] $\in$ {{}}
\end{logs}
The indicated line is inside the function \verb+msgsnd+:
\begin{lstlisting}
  int result = Frama_C_queue_send(msgqid, mess, size);
\end{lstlisting}
Here we have \verb+mess=buf+, \verb+size=2+, and the program is defined only
if  \verb+buf[0..(2-1)]+ is a valid array slice. This
is indeed false in the modified program, as \verb+buf+ has size 1. In
this case, the \Eva analysis is sure that the range is always
invalid, as there is no approximation on either \verb+buf+ or
\verb+size+, and the read fails. This can be verified with the
empty message content in the second line of the log.

\subsubsection{Small buffer during reception}
\nopagebreak[4]
\begin{logs}
../share/mthread_queue.c:16:[kernel] warning: out of bounds write.
          assert \valid(mess+(0..size-1));
../share/mthread_queue.c:16:[kernel] warning: all target addresses were
        invalid. This path is assumed to be dead.
philo.c:54:[mt] warning: Found message of length 2, which is too long for
        buffer 'mess'. Execution will continue without those messages.
       (Ignore "This path is assumed to be dead message if any".)
philo.c:54:[mt] Receiving message on &queue, max size 2, stored in &end.
        No valid value to receive.
\end{logs}
Again, the \Eva analysis detects that we are accessing past the end
of an array.  The warning ``This path is assumed to be dead'' is
actually not really relevant here, and should be ignored. Next,
\mthread adds a more precise warning about which buffer is too small,
and warns that messages of length 2 are too long. This means that any
message of at least that size will be ignored by
\verb+Frama_C_queue_receive+.  Since all messages are of size 2,
there is nothing valid to receive (hence the last line of the log),
and \mthread instructs the \Eva analysis to stop when evaluating the
call to \verb+Frama_C_queue_receive+.


\item[Too many objects.]
  By default, \mthread allows the creation of 32 threads, mutexes or queues,
  with different counters for each kind of object. This value
  is hard-coded in \verb+mthread.h+, in order to have valid~\C.
  If \mthread detects that a program wants to allocate more than this
  number of objects, it issues a warning.
  \begin{logs}
philo.c:51:[mt] warning: During thread creation. Too many thread ids,
         unable to register another one. Try to increase MTHREAD_NUMBER_IDS
         above 32 in the preprocessing directive. Ignoring.
  \end{logs}
  As hinted by the message, the number of possible distinct \mthread
  objects is defined by the \C
  macro \verb+MTHREAD_NUMBER_IDS+. Thus, it suffices to increase its
  value in the preprocessing directive \verb+-cpp-command+, \eg. by
  adding \verb+-DMTHREAD_NUMBER_IDS=40+ for \gcc or \tool{cpp}.


\item[Unrecognized id.] The ids returned by \mthread for threads,
  mutexes and queues are \C \verb+int+s, unless they are cast to
  another type by the program itself. If the code does strange things
  with those \verb+int+s, \eg. incrementing them, it can build precise
  but incorrect ids. \mthread will then fail with a message similar to
  the one below.
  \begin{logs}
philo.c:29:[mt] warning: During mutex lock. Id 13 for mutexes does not
                exists (incrementation inside program?). Ignoring.
  \end{logs}

\item[Uninitialized concurrency structures.]

Primitives receiving an id as argument can be passed the value 0. This
typically corresponds to non-initialized mutexes, queues \etc Either
this is a mistake in the code (the programmer forgot the initialization),
or the initialization will be done later, by another thread, and the
warning should disappear in later iterations of the analysis.

\begin{logs}
philo.c:38:[mt] warning: Trying to unlock uninitialized mutex. Ignoring
\end{logs}

\end{description}




\chapter{Reading \mthread results}
\label{chap:results}

This chapter explains how to interpret the results output by
\mthread, on the  philosophers example.


\section{Reading the results of the philosophers examples}
\label{sec:results-philo}

Running \mthread on \verb+philo.c+ goes smoothly once a proper
\verb+slevel+ (of at least 5) is used. No warning is emitted
during the analysis.
%
\mthread reports it
stops after 4 iterations, having reached the fixpoint.  However,
not all threads are executed at each iteration. For example,
\mthread detects it would learn nothing by analyzing the thread \verb+main+
during its second step,
and thus skips this analysis. If we read more finely
the output, for example by setting \verb+-mt-verbose 2+, the iteration
structure looks like this:
\begin{description}
\item[Initial run of the main thread]
  This analysis detects the five secondary threads. Receiving
  a message on \verb+&queue+ fails. No potential shared zone
  is detected --- quite logically, as only one thread was running.

\item[First iteration]
  The five secondary threads are executed. Messages sent on
  \verb+&queue+ are memorized for an eventual use in another thread.
  %% CHECKRELEASE
  % Taken from examples/mthread/philo.mtverbose2.log
  \begin{logs}
[mt] philo.c:25:
  Sending message on queue, content [0] $\in$ {1}
                                    [1] $\in$ UNINITIALIZED
  \end{logs}

\item[Second iteration]

The main thread is recomputed, because \mthread detects that some
messages can be received on \verb+&queue+.
%% CHECKRELEASE
% Taken from examples/mthread/philo.mtverbose2.log
\begin{logs}
[mt] *** Computing thread <main>, iteration 2 (new message received)
\end{logs}
During this iteration, the call to \verb+msgrcv+ succeeds, and the
value \verb+end[0]+ becomes possibly non-null. As a result, a new
shared memory zone is detected, the variable \verb+end2+.
%% CHECKRELEASE
% Taken from examples/mthread/philo.mtverbose2.log
\begin{logs}
[mt] Concurrent imprecise accesses have changed: before
    \nothing
  vs.
    end2
\end{logs}

\item[Third iteration]

  All threads are recomputed because we want to monitor the accesses to
the potential shared variable \verb+end2+:
  %% CHECKRELEASE
  % Taken from examples/mthread/philo.mtverbose2.log
  \begin{logs}
[mt] *** Computing thread jobs[0], iteration 3
    (potential shared vars changed, interferences changed)
  \end{logs}

At the end of the iteration, \verb+end2+ is detected as being a (really) shared
zone, not just a potential one:
%% CHECKRELEASE
% Taken from examples/mthread/philo.mtverbose2.log
\begin{logs}
[mt] Shared memory: end2
[mt] Concurrent precise var accesses have changed: before
    \nothing
  vs.
    end2
\end{logs}
\mthread also detects that \verb+end2+ is not protected in a coherent way,
\ie. that there might be a race condition on it.
%% CHECKRELEASE
% Taken from examples/mthread/philo.mtverbose2.log
\begin{logs}
[mt] Possible read/write data races:
  end2:
    read by <main> at philo.c:53, unprotected
    write by jobs[0] at philo.c:24, protected by locks[1] locks[4]
    write by jobs[2] at philo.c:24, protected by locks[1] locks[3]
    write by jobs[3] at philo.c:24, protected by locks[2] locks[4]
    write by jobs[4] at philo.c:24, protected by locks[0] locks[3]
[mt] Mutexes for concurrent accesses:
  end2  write protected by (?)locks[0] (?)locks[1] (?)locks[2] (?)locks[3]
        (?)locks[4], read unprotected
\end{logs}
\mthread does not report any new potential shared variable however, which is
coherent with the program.


\item[Fourth iteration]

During this iteration, the thread \verb+main+ is recomputed. Indeed, new
possible values for \verb+end2+ (coming from the other threads), have been
found in iteration 3.
%% CHECKRELEASE
% Taken from examples/mthread/philo.mtverbose2.log
\begin{logs}
[mt] *** Computing thread <main>, iteration 4
    (shared vars values changed, interferences changed)
\end{logs}
During this iteration, the return statement of the \verb+main+ function
becomes reachable.

As this state of the analysis, there is no reason to recompute
any of the threads, and \mthread detects that a fixpoint is reached.
%% CHECKRELEASE
% Taken from examples/mthread/philo.mtverbose2.log
\begin{logs}
[mt] ******* Analysis performed, 4 iterations
\end{logs}
\end{description}
%
Not all the logs given above are available with the default verbosity
level of~1. Indeed, they are not important to understand the
\emph{results} of the analysis, only the way it proceeded.

\medskip

Let us point out a few more information. For example, the information
on the mutexes protecting the accesses to \verb+end2+ are two-fold.
First, we have an exhaustive account, with all accesses by each thread;
each  access is listed together with the mutex
contexts at those points of the analyzes. In this example, the
information is as precise as possible. Second, we have a summary, that
aggregates the exhaustive listing.
%% CHECKRELEASE
% Taken from examples/mthread/philo.mtverbose2.log
\begin{logs}
[mt] Mutexes for concurrent accesses:
  end2  write protected by (?)locks[0] (?)locks[1] (?)locks[2] (?)locks[3]
        (?)locks[4], read unprotected
\end{logs}
This shows that \verb+end2+ is not protected at all when it is read.
Conversely, it is protected by various mutexes when it is written, but
never in a consistent way: there is always a \verb+(?)+ in front of
the mutex name, indicating that in at least one case, the mutex
was not locked. This indicates possible race conditions both when reading
and writing \verb+end2+, which is indeed the case in the program.

\smallskip

Finally, let us discuss the values of the messages sent and received
on \verb+&queue+. We reprint some relevant messages below.
%% CHECKRELEASE
% Taken from examples/mthread/philo.mtverbose2.log
\begin{logs}
[mt] philo.c:25:
  Sending message on queue, content [0] $\in$ {1}
                                    [1] $\in$ UNINITIALIZED
\end{logs}
\begin{logs}
[mt] philo.c:54:
  Receiving message on queue, max size 2, stored in &end. Possible  values:
  From thread jobs[0]: [0] $\in$ {1}
                       [1] $\in$ UNINITIALIZED
  From thread jobs[2]: [0] $\in$ {3}
                       [1] $\in$ UNINITIALIZED
  From thread jobs[3]: [0] $\in$ {4}
                       [1] $\in$ UNINITIALIZED
  From thread jobs[4]: [0] $\in$ {5}
                       [1] $\in$ UNINITIALIZED
\end{logs}
\mthread is quite accommodating about the content of the message,
and tolerates the fact that a part of the source buffer is uninitialized.
Inspecting the value of \verb+end+\footnote{For example using \FramaC's GUI}
after line 54 of \verb+philo.c+ reveals that the possible values are
\begin{logs}
end[0] $\in$ {1; 3; 4; 5}
   [1] $\in$ UNINITIALIZED
\end{logs}
This is also the most precise approximation possible.


\section{\html}
\label{sec:html}

\mthread \html output, triggered by option \verb+-mt-extract html+,
produces a summary of the concurrent program, as well as control-flow
graphs of each thread as analyzed by \mthread.
%
Let us start by the first representation extracted from \mthread's internal
control-flow-graph model: a set of \html pages. This allows easy browsing through the
various information computed by \mthread.

For our simple dining philosophers' example, these pages can be produced in the
directory {\tt html\_summary} by typing :
%% CHECKRELEASE
\begin{shell}
% frama-c -mthread -mt-threads-lib pthreads -mt-extract html \
  $MTSHARE/mthread_queue.c \
  -eva-slevel 256 philo.c
\end{shell}

\begin{figure}[htb]
  \centering
  \includegraphics[scale=.45]{\mthreadtexdir/html_summary.png}
%  \medskip
%  \hrule
  \caption{\html summary of dining philosophers}
  \label{fig:html_summary}
\end{figure}

An \html summary of the code (\myfig{html_summary}) is displayed at {\tt
  html\_summary/index.html},  providing information about thread creations,
lock and unlock directives as well  as message queue uses. There also are links
to the various threads encountered in the program. Clicking on one of those
links leads to a summary concerning the given thread.
This thread-focused summary (\myfig{thread_summary}) shows the concurrent
control-flow graph of the thread (\mysec{mthread-cfg}).



%% CHECKRELEASE
\begin{figure}[htb]
  \centering
  % You can "Inspect" the page to edit the svg in the browser and increase the
  % scale of the SVG to be able to read something in the screenshot.
  \includegraphics[scale=.35]{\mthreadtexdir/thread_summary.png}
%  \medskip
%  \hrule
  \caption{Excerpt from the first philosopher's control-flow graph}
  \label{fig:thread_summary}
\end{figure}

The precision and details
shown on this graph can be controlled by \mthread's options detailed in
\ref{sec:mthread-cfg} and \myappdx{options}.
Links to all the other threads are provided but one important
feature here is that the control-flow graph is clickable. A click on the
{\tt Call aux} node yields the source code behind this node
(\myfig{thread_code}). Most of the expressions in the source
code page are themselves clickable, for example to navigate from
function to function.

%% CHECKRELEASE
\begin{figure}[htb]
  \centering
  \includegraphics[scale=.35]{\mthreadtexdir/thread_code.png}
%  \medskip
%  \hrule
  \caption{Excerpt from the philosophers' source code}
  \label{fig:thread_code}
\end{figure}


\chapter{Command-line options}
\label{sec:options}

This section describes the list of options available to finely
tune the behavior of \mthread. They can also be manually found
using {\tt frama-c -mt-help}. Some experimental options are intentionally
left undocumented.

\paragraph{Basics.}
As a reminder, the generic options for \mthread are the following:
\begin{description}
\item[{\tt -mthread}] This enables the \mthread plug-in. This option is
  mandatory for any use of \mthread and launches the \mthread analysis.

\item[{\tt -mt-threads-lib {builtins-only|pthreads}}] Specify which threading
  library is used. If the user provides custom stubs then \verb+builtins-only+
  should be used to only bring into scope \mthread's built-in primitives. This
  option can be used without \verb+-mthread+ to add the stubs to the parsing
  stage without running an analysis in a two-steps process:
  %% CHECKRELEASE
  \begin{shell}
  % frama-c -mt-threads-lib pthreads \
    ccfg.c \
    -save session.sav
  % frama-c -load session.sav -mthread -eva-verbose 0
  \end{shell}

\item[{\tt -mt-verbose n}] Change the verbosity of \mthread. Default is
  1. Any value strictly above 1 will show the internal state of the
  analysis at the end of each iteration.

\item[{\tt -mt-help}] Display a short summary of all the \mthread
  options available.

  (The options for \FramaC in general can be obtained through
  \verb+frama-c -kernel-help+, while those for the \Eva analysis are
  invoked by \verb+frama-c -eva-help+.)


\end{description}

\paragraph{External outputs.}
The \mthread results printed as \html for further study.
\begin{description}
\item[{\tt -mt-extract html}] Extracts a partial version of the
  results found by \mthread as HTML.  All results can be
  browsed\footnote{A navigator with support for \tool{SVG} files
    is required to display the control-flow graphs.}
  starting from the file \verb+./html_summary/index.html+.
\end{description}

\paragraph{Control-flow graph options.}

The options below control how the concurrent control-flow graphs are displayed
and simplified.
\begin{description}

\item[{\tt -mt-return-edges}] Link nodes for function calls to
  their corresponding {\tt return} nodes. This makes it easier
  to see nested calls of big functions. (Set by default)

\item[{\tt -mt-non-shared-accesses}] Do not remove nodes corresponding
  to accesses to false shared accesses (\mysec{shared-vars}). Not set
  by default; if the option is set, the accesses are shown in white
  in the control-flow graph.

\item[{\tt -mt-non-concurrent-accesses}] Do not remove nodes corresponding
  to accesses to shared accesses that occur in a non-concurrent context
  (\mysec{shared-vars}). Set by default, those accesses are shown in
  green in the control-flow graph.

\item[{\tt -mt-full-cfg}] Do not simplify the bodies of functions that
  contain multithreaded events. All the statements of those functions
  will be reflected in the control-flow graphs, which can result in very
  big graphs:
  use this option with caution. Calls to functions that do not contain
  multithreaded events are however never inlined in the control-flow graphs.
  Not set by default.


\end{description}


\paragraph{Debug  options.}

Those debug options are not intended for general use, but can
sometimes be useful to diagnose a strange behavior of \mthread. Other
debug options are unintentionally not described.

\begin{description}
\item[{\tt -mt-stop-after <i>}] Instructs \mthread to only perform
  at most {\tt i} iterations of the analysis. If the analysis has
  not converged by then, it is stopped, and the remaining steps
  to perform are shown on the log.

\end{description}


\appendix

\chapter{\mthread functions for stubbing}
\label{sec:stubbed-functions}
\label{appendix:stubbed-functions}

This appendix details the concurrent functions \mthread is able to
detect and handle.  Their prototypes can be found in the file
\verb+$MTSHARE\mthread.h+.

\begin{description}
\item[Thread-related primitives] \

  \begin{itemize}
  \item Thread creation, through function \verb+Frama_C_thread_create+
  \item Thread start, through function \verb+Frama_C_thread_start+
  \item Thread immediate exit, through function \verb+Frama_C_thread_exit+
  \item Current thread id, through function \verb+Frama_C_thread_id+
  \item Thread suspension, through function \verb+Frama_C_thread_suspend+
  \item Thread canceling, through function \verb+Frama_C_thread_cancel+

    \it (This functions currently cancels the thread regardless of any
    potential cancellability state notion, such as the one available in
    \pthreads.)
  \end{itemize}


\item[Mutex-related primitives] \

  \begin{itemize}
  \item Mutex initializing, through function \verb+Frama_C_mutex_init+
  \item Mutex locking, through function \verb+Frama_C_mutex_lock+
  \item Mutex release, through function \verb+Frama_C_mutex_unlock+
  \end{itemize}

\item[Queue-related primitives] \

  \begin{itemize}
  \item Queue initializing, through function \verb+Frama_C_queue_init+
  \item Message sending, through function \verb+Frama_C_queue_send+
  \item Message reception, through function \verb+Frama_C_queue_receive+
  \end{itemize}

\item[Miscellaneous functions] \

  \begin{itemize}
  \item Logging, through function \verb+Frama_C_mthread_show+

    This function takes as first argument a constant string will be
    used as a message, and a number of \C values that will be printed
    after the message. It can be used to show in the control-flow graph any information
    relevant to the analysis, and does not modify the memory state at all.

  \item Forcing synchronization of unprotected shared values, through
    the use of the function \verb+Frama_C_mthread_sync+.

  \end{itemize}


\end{description}

More involved concurrency primitives, such as spinlocks \etc. are not
currently supported. They may be added to \mthread later.


\bibliographystyle{alpha}
\bibliography{biblio}

\end{document}
